\documentclass[10pt,a4paper]{amsart}
\usepackage[margin=19mm]{geometry}
\usepackage{amsmath,amssymb,mathtools}
\usepackage[scr=rsfs,cal=euler]{mathalfa}
\usepackage{xfrac}
\usepackage[unicode,hidelinks,bookmarks=false]{hyperref}
\newcommand{\ie}{i.e.\ }
\newcommand{\cf}{cf.\ }
\newcommand{\resp}{respectively\ }
\newcommand{\Subsection}{\subsection}
\providecommand{\nobreakdash}{\nobreak\hbox{-}}
\setlength{\emergencystretch}{2em}
\multlinegap0pt
\usepackage{amssymb}
\usepackage{tikz}
\usepackage[shortlabels]{enumitem}
\usepackage{xcolor}
\newtheorem{theorem}{Theorem}
\title{Transfer of difference structures: a new semidirect product framework}
\author{Sophie Huczynska, Struan McCartney, Carys Williams}
\date{Source: arXiv:2609.18897v1 (CC BY 4.0)}
\begin{document}
\maketitle

\noindent\emph{Source and licence.} Transfer of difference structures: a new semidirect product framework. Version arXiv:2609.18897v1 (CC BY 4.0), distributed by arXiv under the Creative Commons Attribution 4.0 International licence, http://creativecommons.org/licenses/by/4.0/, as recorded in the arXiv licence metadata for this exact version and shown on the abstract page https://arxiv.org/abs/2609.18897v1. Full licence notice: \texttt{licenses/arxiv-2609.18897-CC-BY-4.0.txt}.\par\smallskip

\noindent\textit{Editorial scope.} Counted unit: the article's theorem thm:noncyclic (source0.tex lines 682--691). The article's complete original proof of it is delivered with it (source0.tex lines 693--830, one \texttt{proof} environment of 138 lines). No mathematics is written, repaired or re-derived: the statement and the proof are the article's own lines, in the article's own notation, and no retained line is substituted or re-flowed.  No further source-local context is needed: every cross-reference of every retained line resolves inside the two delivered spans.  Nothing was omitted from any retained span, so the package carries no omission record.  No retained line contains a citation, so the wrapper carries no bibliography and no bibliography entry.  No retained line contains a figure environment, an image-inclusion command or drawing code, so no image is delivered and the package declares no asset dependency.  Editorial formatting only, in addition: an AMS \texttt{amsart} preamble that restates the article's own declarations and macros used by the retained lines, namely the environments \texttt{theorem} on separate counters, together with the standard packages those lines need.  The retained environments print in this document as Theorem 1 in order of appearance, whereas the article prints the same items with the numbering of its own full document.  The complete native source of the article as distributed in the arXiv e-print for this exact version is bundled unchanged as \texttt{sources/210734/source0.tex}.
\begin{theorem}\label{thm:noncyclic}
Let $k \in \mathbb{N}$. Define the following subsets of $\mathbb{Z}_{\frac{3(2k+1)^2+1}{2}}\times\mathbb{Z}_2$:
\begin{itemize}
\item $C_0=\bigcup_{i=0}^{2k}\{(-k+i,0)\}$
\item $C_1=\bigcup_{i=0}^{2k}\{(k+1+(3k+2)i,i)\}$
\item $C_2=\bigcup_{i=0}^{2k}\{(2k+1+(3k+1)i,i+1)\}$
\end{itemize}
where the first component is taken modulo $\frac{3(2k+1)^2+1}{2}$ and the second component is taken modulo $2$.\\
Then $(C_0,C_1,C_2)$ form a $(3(2k+1)^2+1,3,2k+1,1)$-CEDF in the non-cyclic abelian group $\mathbb{Z}_{\frac{3(2k+1)^2+1}{2}}\times\mathbb{Z}_2$ and each $C_i$ ($0 \leq i \leq 2$) is a symmetric set in this group.
\end{theorem}

\begin{proof}
Note $(3(2k+1)^2+1)/2=6k^2+6k+2$ is even and hence $G=\mathbb{Z}_{\frac{3(2k+1)^2+1}{2}}\times\mathbb{Z}_2$ is not a cyclic group.  
Throughout this proof, for simplicity, we will denote $D(X,Y; G)$ by $D(X,Y)$. We will show that $D(C_1,C_0) \cup D(C_2,C_1) \cup D(C_0,C_2)=G \setminus \{0\}$; note that disjointness of the three sets then follows from the fact that $0$ does not occur in the multiset of differences.  Since the difference multiset contains $3(2k+1)^2$ elements by construction, it will suffice to show that each non-identity group element occurs at least once.\\
For $a, b \in \mathbb{Z}_{\frac{3l^2+1}{2}}$ and $x \in \mathbb{Z}_2$, we will use the interval notation: we denote the set $\{(a,x), (a+1,x), \ldots, (b,x)\}$ by $[a,b]\times \{x\}$. \\
First, we determine the difference multisets. 
\begin{align*}
D  (C_1,C_0) &= \bigcup_{i=0}^{2k}\bigcup_{j=0}^{2k}\{(k+1+(3k+2)i-(-k+j),i)\}\\
&= \bigcup_{i=0}^{2k}\bigcup_{j=0}^{2k}\{((3k+2)i+2k+1-j,i)\}\\
&= \bigcup_{i=0}^{2k}[(3k+2)i+1,(3k+2)i+2k+1]\times \{i\}.\\
\end{align*}
We may view these differences as consisting of $2k$ length-$2k$ ``runs" of consecutive elements in the first coordinate, which occur horizontally in the difference table, indexed by $i$. Each run has fixed second coordinate, and the parity of the second coordinate alternates as $i$ takes values from $0$ to $2k$.

Next we have:
\begin{align*}
D(C_2,C_1) &= \bigcup_{i=0}^{2k}\bigcup_{j=0}^{2k}\{(2k+1+(3k+1)i-(k+1+(3k+2)j),i-j+1)\}\\
&= \bigcup_{i=0}^{2k}\bigcup_{j=0}^{2k}\{(3k+1)(i-j+1)-2k-1-j,i-j+1)\}.
\end{align*}
We let $s = i-j+1$: since $i$ and $j$ vary from $0$ to $2k$, $s$ takes values from  $-2k+1$ to $2k+1$.  We change the indices from $i$ and $j$ to $s$ and $j$.  For $-2k+1 \leq s \leq 2k+1$, let  $$J_s=\{j: 0 \leq j \leq 2k \mbox{ and } 0 \leq s+j-1 \leq 2k\}=[0,2k] \cap [1-s,2k+1-s]$$
so that
$$D  (C_2,C_1) = \bigcup_{s=-2k+1}^{2k+1}\bigcup_{j \in J_s} \{((3k+1)s-j-2k-1,s)\}.$$
For $-2k+1 \leq s \leq 0$, we have $J_s=[1-s,2k]$, while for $1 \leq s \leq 2k+1$ we have $J_s=[0,2k+1-s].$  Thus we have:
\[ D  (C_2,C_1) = \bigcup_{s=-2k+1}^{0}
\bigcup_{j = 1-s}^{2k} \{((3k+1)s-j-2k-1,s)\}
\cup\bigcup_{s=1}^{2k+1}\bigcup_{j =0}^{2k+1-s} 
\{((3k+1)s-j-2k-1,s)\}\]
\[= \bigcup_{s=-2k+1}^{0}
[(3k+1)s-4k-1,(3k+2)s-2k-2]\times \{s\}\]\[\cup\bigcup_{s=1}^{2k+1}
[(3k+2)s-4k-2,(3k+1)s-2k-1]\times \{s\}.\]
We may view these differences as ``runs" of consecutive elements (in the first coordinate) occurring diagonally in the difference table The set $J_s$ allows for the adjustment of the length of the differences depending on which diagonal we are considering. Again, we have alternation between $0$ and $1$ in the second coordinate for each run.

Finally:
\begin{align*}
D(C_0,C_2) &=\bigcup_{i=0}^{2k}\bigcup_{j=0}^{2k}\{(-k+i-(2k+1+(3k+1)j),j+1)\}\\
&=\bigcup_{j=0}^{2k}[-(3k+1)j-3k-1,-(3k+1)j-k-1]\times\{j+1\}.
\end{align*}
These differences we may view as ``runs" of consecutive elements (in the first coordinate) occurring vertically in the difference table, indexed by $j+1$, with alternation in the second coordinate. Re-indexing to replace $j$ by $2k-j$, and using the fact that $6k^2+6k+2 \equiv 0 \mod 3(2k + 1)^2 + 1)/2$ and $j \equiv (2k-j) \mod 2$, we have that
\begin{align*}
D(C_0,C_2)=\bigcup_{j=0}^{2k}[(3k+1)j+k+1,(3k+1)(j+1)]\times\{j+1\}.
\end{align*}

Having obtained expressions for each $D(C_{i+1 \mod 3},C_{i})$ ($i \in \{0,1,2\}$), we next check that the difference multiset comprises one occurrence of each non-identity element in $\mathbb{Z}_{\frac{3(2k+1)^2+1}{2}} \times \{0\}$.  We consider the differences corresponding to the following indices, where $ 0 \leq t \leq k-1$.  
\begin{itemize}
\item $i=2t$ in $D(C_1,C_0)$: the difference multiset obtained is
\begin{align*}
&= \bigcup_{t=0}^{k-1}[(3k+2)(2t)+1,(3k+2)(2t)+2k+1]\times \{0\}\\
\end{align*}
\item $s=2(t+1)$ in $D(C_2,C_1)$: the difference multiset obtained is
\begin{flalign*}
&=\bigcup_{t=0}^{k-1}
[(3k+2)(2(t+1))-4k-2,(3k+1)(2(t+1))-2k-1]\times \{0\}\\
&=\bigcup_{t=0}^{k-1} 
[(3k+2)(2t)+2k+2,(3k+1)(2t+1)+k]\times \{0\}\\
\end{flalign*}
\item $j =2t+1$ in $D(C_0,C_2)$: the difference multiset obtained is
\begin{align*}
&=\bigcup_{t=0}^{k-1}[(3k+1)(2t+1)+k+1,(3k+1)(2(t+1))]\times\{0\}
\end{align*}
\item $s=-2t$ in $D(C_2,C_1)$ (here $-2k+1 \leq s \leq 0$):
the difference multiset obtained is
\begin{flalign*}
&= \bigcup_{t=0}^{k-1}
[(3k+1)(-2t)-4k-1,(3k+2)(-2t)-2k-2]\times \{0\}\\
&= \bigcup_{t=0}^{k-1}
[(3k+1)(-2(k-1-t))-4k-1+6k^2+6k+2,\\
&(3k+2)(-2(k-1-t))-2k-2+6k^2+6k+2]\times \{0\}\\
&= \bigcup_{t=0}^{k-1}
[(3k+1)(2(t+1))+1,(3k+2)(2(t+1)))]\times \{0\}\\
\end{flalign*}
Here we have reindexed to replace $t$ by $k-1-t$ and used the fact that $6k^2+6k+2 \equiv 0 \mod \frac{3(2k+1)^2+1}{2}$.
\end{itemize}
The union of these is $[(3k+2)2t+1,(3k+2)2(t+1)]\times\{0\}$.
As $t$ ranges from $ 0 \text{ to } k-1$ we obtain one copy of:
\[[1,6k^2+4k]\times\{0\}\]
Finally, we take $i=2k$ in $D(C_1,C_0)$.  This contributes the set of differences 
$$[6k^2+4k+1,6k^2+6k+1]\times\{0\}.$$

Now we consider the elements of $\mathbb{Z}_{\frac{3(2k+1)^2+1}{2}}\times\{1\}$. Again we consider the differences corresponding to the following indices, where $ 0 \leq t \leq k-1$.  
\begin{itemize}
\item $j=2t$ in $D(C_0,C_2)$: the difference multiset obtained is
\begin{align*}
&=\bigcup_{t=0}^{k-1}[(3k+1)(2t)+k+1,(3k+1)(2t+1)]\times\{1\}
\end{align*}
\item $s=-2(t+1)+1$ in $D(C_2,C_1)$ (here $-2k+1 \leq s \leq 0$): the difference multiset obtained is
\begin{flalign*}
&= \bigcup_{t=0}^{k-1}
[(3k+1)(-2(t+1)+1)-4k-1,(3k+2)(-2(t+1)+1)-2k-2]\times \{1\}\\
&= \bigcup_{t=0}^{k-1}
[(3k+1)(-2(k-1-t+1)+1)-4k-1+6k^2+6k+2,\\
&(3k+2)(-2(k-1-t+1)+1)-2k-2+6k^2+6k+2]\times \{1\}\\
&= \bigcup_{t=0}^{k-1}
[(3k+1)(2t+1)+1,(3k+2)(2t+1)]\times \{1\}\\
\end{flalign*}
Here we have reindexed to replace $t$ by $k-1-t$ and used the fact that $6k^2+6k+2 \equiv 0 \mod \frac{3(2k+1)^2+1}{2}$.
\item $i =2t+1$ in $D(C_1,C_0)$: the difference multiset obtained is
\begin{align*}
&=\bigcup_{t=0}^{k-1}[(3k+2)(2t+1)+1,(3k+2)(2t+1)+2k+1]\times\{1\}
\end{align*}
\item $s=2(t+1)+1$ in $D(C_2,C_1)$: the difference multiset obtained is 
\begin{flalign*}
&=\bigcup_{t=0}^{k-1}
[(3k+2)(2(t+1)+1)-4k-2,(3k+1)(2(t+1)+1)-2k-1]\times \{1\}\\
&=\bigcup_{t=0}^{k-1} 
[(3k+2)(2t+1)+2k+2,(3k+1)(2(t+1))+k]\times \{1\}\\
\end{flalign*}
\end{itemize}
The union of these is $[(3k+2)2t+k+1,(3k+1)2(t+1)+k]\times\{1\}$.
As $t$ ranges from $ 0 \text{ to } k-1$ we obtain one copy of:
\[[k+1,6k^2+3k]\times\{1\}\]

Finally, we take $s=1$ in $D(C_2,C_1)$ and $j=2k$ in $D(C_0,C_2)$.  These contribute the sets of differences: 
$$([6k^2+5k+2,6k^2+6k+1]\cup[0,k])\times\{1\}.$$

and
$$[6k^2+3k+1,6k^2+5k+1]\times\{1\},$$

We have shown that $D(C_1,C_0) \cup D(C_2,C_1) \cup D(C_0,C_2)$ comprises one copy of each non-identity element in $\mathbb{Z}_{\frac{3(2k+1)^2+1}{2}}\times\mathbb{Z}_2$, and so $(C_0,C_1,C_2)$ form a $(3(2k+1)^2+1,3,2k+1,1)$-CEDF in $\mathbb{Z}_{\frac{3(2k+1)^2+1}{2}}\times\mathbb{Z}_2$.

To show that each set is symmetric, we re-index. Letting $j =i-k$, we have:
\begin{itemize}
\item $C_0=\bigcup_{j=-k}^{k}\{(j,0)\}$
\item $C_1=\bigcup_{j=-k}^{k}\{(3k^2+3k+1+(3k+2)j,j+k)\}$
\item $C_2=\bigcup_{j=-k}^{k}\{(3k^2+3k+1+(3k+2)j,j+k+1)\}$
\end{itemize}
Let $C_{0,j}=\{(j,0)\}$, $C_{1,j}=\{(3k^2+3k+1+(3k+2)j,j+k)\}$, $C_{2,j}=\{(3k^2+3k+1+(3k+2)j,j+k+1)\}$.  With this notation, the sets may be re-written as:
\begin{itemize}
\item $C_0= C_{0,0} \cup \bigcup_{j=1}^{k}\{C_{0,j},C_{0,-j}\}$
\item $C_1= C_{1,0} \cup \bigcup_{j=1}^{k}\{C_{1,j},C_{1,-j}\}$
\item $C_2= C_{2,0} \cup \bigcup_{j=1}^{k}\{C_{2,j},C_{2,-j}\}$
\end{itemize}
It can be verified that $C_{i,-j}=-C_{i,j}$ for all $i \in \{0,1,2\}$ and all $0 \leq j \leq k$. We demonstrate this for $i=1$; the other cases follow similarly. 
\begin{align*}
C_{1,-j}&=\{(3k^2+3k+1+(3k+2)(-j),-j+k)\}\\
    &=\{(3k^2+3k+1-(3k+2)j-(6k^2+6k+2),-j-k)\}\\
    &=\{-(3k^2+3k+1+(3k+2)j,j+k)\} \\
    &= -C_{1,j}
\end{align*}
Hence $C_{i}=-C_i$ for $i \in \{0,1,2\}$ and so the sets are symmetric.
\end{proof}

\end{document}
